\section{Integral Lie algebras and the special identity}
\label{sec:algebra}

Throughout, a Lie algebra is an ordinary Lie algebra: its bracket is
alternating, including in characteristic two.  Put
\(R=\Z_{(2)}=\{a/b\in\Q:b\text{ is odd}\}\), with residue field
\(\F=R/2R\).  We will pass from the rational solution space to an
embedded reduction of the same dimension.  For this passage we need
torsion-free integral targets for the defining equations, together with
an additional identity proved over \(\Q\) before reduction.

\begin{lemma}\label{lem:low-weights}
The solution spaces \(W_1\) and \(W_2\) are zero.
\end{lemma}
\begin{proof}
For \(\psi(x,y)=\alpha x+\beta y\), the left side minus the right side of
the pentagon is \(-\alpha t_{12}-\beta t_{34}\).  The degree-one braid
generators are linearly independent, so \(\alpha=\beta=0\).
The degree-two free Lie algebra is spanned by \([x,y]\), and the
three-term expression for this polynomial is \(3[x,y]\).  Thus it too
has no nonzero rational solution.
\end{proof}

\begin{samepage}
\begin{lemma}[Ordinary Lie lattices]\label{lem:lie-lattice}
Let \(X\) be a finite alphabet, and let \(K(X)\) be the \(R\)-span of
all bracket monomials in the free associative algebra
\(R\langle X\rangle\).  It is the free ordinary Lie algebra over \(R\).
Each multihomogeneous piece is finite free, and its rationalization and
reduction are, respectively,
\[
 K(X)\otimes_R\Q=\Lie_{\Q}\langle X\rangle,
 \qquad
 K(X)/2K(X)=\Lie_{\F}\langle X\rangle
       \ \subset\ \F\langle X\rangle.
\]
In particular, this reduction is injective in associative words, and
free ordinary Lie dimensions in every multidegree agree over \(\Q\)
and \(\F\).  The same statements hold for an alphabet with positive
integer weights and finitely many letters in each bounded weight.
\end{lemma}
\end{samepage}
\begin{proof}
We recall the integral free-Lie and PBW facts being used.  The
Hall--Lyndon construction gives a basis of the free Lie algebra over
\(\Z\): iterated use of alternation and Jacobi expresses bracket
monomials in the standard Lyndon brackets, and the associative
expansion of each such bracket has its Lyndon word as leading word,
with coefficient one.  Distinct leading words give independence over
every coefficient ring.  Thus base change to \(R\) identifies the
abstract free Lie algebra with precisely the displayed bracket span.
PBW applies to this free \(R\)-module; in particular its universal
enveloping algebra is \(R\langle X\rangle\).  These are the usual
integral versions of the free-Lie basis and PBW theorems
\cite[Theorems~0.2, 0.5, 4.9 and~5.1]{Reutenauer};
the integral Lyndon and PBW bases are also recalled in
\cite[Section~2]{MelanconReutenauer}.

The same Lyndon brackets form bases after base change to $\Q$ and
$\F$, with unchanged multidegrees.  This proves the asserted
injectivity and equality of dimensions.  For a weighted alphabet,
restriction to any bounded weight leaves finitely many letters,
so the same basis argument applies.
\end{proof}

In particular, ordinary Lie membership is a condition stronger than
being primitive in a characteristic-two enveloping algebra.  For
example, \(x^2\) is primitive in \(\F\langle x\rangle\), whereas the
ordinary Lie algebra generated by \(x\) is just \(\F x\).  We will
always retain ordinary Lie membership when using associative words.

For a coefficient ring \(k\in\{R,\Q,\F\}\), write
\(\mathfrak t_q(k)\) for the infinitesimal braid Lie algebra on
\(q\) strands, with the presentation in \eqref{eq:braid-relations}.

\begin{lemma}[The free fiber]\label{lem:free-fiber}
The map \(\mathfrak t_q(k)\longrightarrow\mathfrak t_{q-1}(k)\)
which forgets the last strand is split, and its kernel is the free
ordinary Lie algebra on \(s_i=t_{iq}\), \(1\leq i<q\).  Explicitly,
\begin{equation}\label{eq:free-fiber}
 \mathfrak t_q(k)
  =\Lie_k\langle s_1,\ldots,s_{q-1}\rangle
        \rtimes\mathfrak t_{q-1}(k),
\end{equation}
where a base generator \(t_{ij}\) acts by the derivation
\begin{equation}\label{eq:free-fiber-action}
 D_{ij}(s_i)=[s_i,s_j],\qquad
 D_{ij}(s_j)=[s_j,s_i],\qquad
 D_{ij}(s_h)=0\quad(h\notin\{i,j\}).
\end{equation}
Multiplication gives an isomorphism of \(k\)-modules
\begin{equation}\label{eq:enveloping-fiber}
 k\langle s_1,\ldots,s_{q-1}\rangle
       \otimes_k U(\mathfrak t_{q-1}(k))
       \xrightarrow{\ \sim\ }U(\mathfrak t_q(k)).
\end{equation}
Thus recursive products of free associative fiber words are a basis.
Over \(R\), the Lie algebra and its enveloping algebra are torsion
free, and this normal form commutes with rationalization and reduction.
\end{lemma}
\begin{proof}
Any prescription on free generators extends uniquely to a derivation.
Derivations \(D_{ij}\) and \(D_{hk}\) with disjoint supports commute.
For three distinct indices \(i,j,h\), put
\(S_{ijh}=s_i+s_j+s_h\).  On the free subalgebra on these three letters,
\[
 D_{ij}+D_{ih}+D_{jh}:z\longmapsto[z,S_{ijh}].
\]
Moreover \(D_{ij}(S_{ijh})=0\), so \(D_{ij}\) commutes with this sum.
On every other free generator all these derivations vanish.  It
follows that \([D_{ij},D_{ih}+D_{jh}]=0\) on the whole free Lie
algebra.  Hence \eqref{eq:free-fiber-action} defines an action of the
presented base algebra.

Form its semidirect product with the free Lie algebra on the \(s_i\).
The braid relations involving the last strand say precisely
\[
 [t_{ij},s_h]=0\ (h\notin\{i,j\}),\quad
 [t_{ij},s_i+s_j]=0,\quad
 [s_i,t_{ij}+s_j]=[s_j,t_{ij}+s_i]=0.
\]
They hold by \eqref{eq:free-fiber-action}.  Conversely, these relations
in the presented braid algebra give exactly that action.  The maps
between the presented algebra and the semidirect product, taking each
generator to its namesake, are consequently inverse.  This proves
both the presentation and the assertion about the actual kernel.

Inductively the base is a free \(k\)-module.  Choose Lie bases for
fiber and base, and place the entire fiber basis before the base
basis.  PBW identifies their ordered products with a basis of the
enveloping algebra.  Since the enveloping algebra of a free Lie
algebra is the free associative algebra, this is exactly
\eqref{eq:enveloping-fiber}.  Induction, starting with one strand,
gives the asserted word basis.  Lemma~\ref{lem:lie-lattice} and the
same basis over \(R\), \(\Q\), and \(\F\) give the last assertions.
\end{proof}

We next deduce an additional identity from the defining equations.
The first conclusion below is Drinfeld's special identity
\cite[Proposition~5.7]{Drinfeld}.  We give a direct proof that also
establishes the stated linearization.  We use the convention
\(\ad_a(z)=[a,z]\).

\begin{lemma}[The special identity]\label{lem:special}
Let \(\psi\in W_n\), where \(n>2\).  In the free Lie algebra on
\(a,b\), with \(c=-a-b\), one has
\begin{equation}\label{eq:special}
 [a,\psi(b,a)]+[c,\psi(b,c)]=0.
\end{equation}
In addition, for an extra free letter \(T\),
\begin{equation}\label{eq:special-linearization}
 \partial_2\psi(a,-a)T=0,
\end{equation}
where \(\partial_2\) means the coefficient of a central parameter
\(t\) in \(\psi(a,-a+tT)\).
\end{lemma}
\begin{proof}
All calculations in this proof take place over \(\Q\).  We first
rewrite the pentagon in a free kernel, then linearize at \(a+b+c=0\).
Equality of mixed derivatives will force the possible linearization
coefficient to vanish.

\emph{The free-fiber identity.}
Forget strand \(2\) in \(\mathfrak t_4(\Q)\), and put
\[
 a=t_{12},\qquad b=t_{23},\qquad c=t_{24},\qquad
 u=t_{13},\qquad v=t_{34}.
\]
By Lemma~\ref{lem:free-fiber}, the kernel is the actual free Lie
algebra on \(a,b,c\).  Define
\begin{align*}
 F_a&=\psi(u+b,v)-\psi(u,v),\\
 F_b&=\psi(u+a,v+c)-\psi(u,v),\\
 F_c&=\psi(u,v+b)-\psi(u,v).
\end{align*}
Forgetting the fiber sends each displayed difference to zero, so
\(F_a,F_b,F_c\) are elements of this free kernel.  They are
homogeneous of degree \(n\).

Each of the three expressions before subtracting \(\psi(u,v)\)
commutes with its corresponding letter.  For example,
\([a,u+b]=[a,v]=0\) and
\([b,u+a]=[b,v+c]=0\); the third case is the same braid relation at
\(c\).  The base action fixes \(S=a+b+c\).  Summing the resulting
commutators therefore gives
\[
 [a,F_a]+[b,F_b]+[c,F_c]=0.
\]
Set \(\Delta_a=F_b-F_a\) and \(\Delta_c=F_b-F_c\).  We obtain the
identity in the free fiber
\begin{equation}\label{eq:special-fiber}
 [S,F_b]=[a,\Delta_a]+[c,\Delta_c].
\end{equation}

The pentagon gives
\(\Delta_a=\psi(a,b+c)-\psi(a,b)-\psi(b,v)\).
The element \(v+b+c\) is central in the subalgebra on \(v,b,c\).
A Lie polynomial of degree greater than one is unchanged when a
central element is added to an argument: every term in the difference
contains that element in a bracket.  Thus \(\psi(b,v)=\psi(b,-b-c)\).
Exchanging strands \(1\) and \(4\) in the same calculation and using
antisymmetry gives the second formula below:
\begin{align}
 \Delta_a&=\psi(a,b+c)-\psi(a,b)-\psi(b,-b-c),
       \label{eq:special-delta-a}\\
 \Delta_c&=-\psi(c,a+b)+\psi(c,b)+\psi(b,-b-a).
       \label{eq:special-delta-c}
\end{align}
Both differences vanish when \(S=0\): the two remaining terms cancel
by antisymmetry and \(\psi(z,-z)=0\) in degree greater than one.

The ideal generated by \(S\) is stable under the base action, so the
quotient fiber is \(\Lie_{\Q}\langle a,b\rangle\), with
\(c=-a-b\).  In this quotient the actions of \(u,v\) are
\(\ad_c,\ad_a\), respectively.  For example,
\[
 [u,a]=[a,b]=[c,a],\quad [u,b]=[b,a]=[c,b],\quad [u,c]=0,
\]
and
\([v,a]=0\), \([v,b]=[b,c]=[a,b]\),
\([v,c]=[c,b]=[a,c]\).  It follows that \(\psi(u,v)\) acts on
the quotient fiber by \(\ad_{\psi(c,a)}\).  Since the expression
defining \(F_b+\psi(u,v)\) commutes with \(b\), we have
\[
 [b,F_b+\psi(c,a)]=0
\]
in this free Lie quotient.

The centralizer of \(b\) in the ordinary free Lie algebra on \(a,b\)
is \(\Q b\).  To see this directly, embed it in associative words.
For a homogeneous polynomial \(P=\sum_w p_w w\) commuting with
\(b\), comparison in \(bP=Pb\) gives
\(p_{av}=0\) and
\(p_{b^j av}=p_{b^{j-1}avb}\) for every word \(v\) and \(j>0\).
Iterating shows that every word containing \(a\) has coefficient
zero.  Hence \(P\) is a polynomial in \(b\).  Setting \(a=0\) now shows that a
Lie element of this form belongs to the one-dimensional ordinary
Lie algebra on \(b\).  The homogeneous element at hand has degree
\(n>1\), and is therefore zero.  We have proved an equality of
quotient fiber elements, not just of their induced derivations:
\begin{equation}\label{eq:special-fb}
 F_b\big|_{S=0}=-\psi(c,a).
\end{equation}

\emph{The first variation at \(S=0\).}
We have determined the quotient value of \(F_b\).  We now differentiate
\eqref{eq:special-fiber} transversely to this quotient and compare the
result with the derivative of the expression in \eqref{eq:special}.
All subsequent derivatives mean coefficient
extraction after adjoining free direction letters and commuting
scalar parameters.  Thus
\[
 \partial_1\psi(p,q)T=[t]\psi(p+tT,q),\qquad
 \partial_2\psi(p,q)T=[t]\psi(p,q+tT).
\]
The degree-\(n\) Lie polynomials on \(a,T\) containing exactly one
\(T\) and \(n-1\) copies of \(a\) form the one-dimensional space
spanned by \(\ad_a^{n-1}T\): in a nonzero bracket monomial, each
branch not containing \(T\) must be a single \(a\).  Consequently
there is a scalar \(k\in\Q\), depending only on \(\psi\), such that
\begin{equation}\label{eq:special-k}
 \partial_2\psi(a,-a)T=k\ad_a^{n-1}T.
\end{equation}
Differentiating the identity \(\psi(a,-a)=0\) while varying both
arguments gives
\begin{equation}\label{eq:special-diagonal}
 \partial_1\psi(a,-a)T=\partial_2\psi(a,-a)T.
\end{equation}

Take \(c=S-a-b\), keep \(a,b\) fixed, and extract the coefficient
linear in \(S=tT\).  In these same coordinates the two differences
are
\begin{align*}
 \Delta_a&=\psi(a,S-a)-\psi(a,b)-\psi(b,a-S),\\
 \Delta_c&=-\psi(S-a-b,a+b)+\psi(S-a-b,b)+\psi(b,-b-a).
\end{align*}
Writing \(c=-a-b\) after differentiation, their derivatives at zero
are exactly
\begin{align*}
 d_S\Delta_a(T)
   &=k\ad_a^{n-1}T+\partial_2\psi(b,a)T,\\
 d_S\Delta_c(T)
   &=-\partial_1\psi(c,-c)T+\partial_1\psi(c,b)T\\
   &=-k\ad_c^{n-1}T-\partial_2\psi(b,c)T.
\end{align*}
The last equality uses \eqref{eq:special-diagonal} and
antisymmetry.  In differentiating \eqref{eq:special-fiber}, the left
side becomes \([T,F_b|_{S=0}]\): its other product-rule term has
outside letter \(S=0\).  On the right the outside \(a\) is fixed,
and the derivative of the outside \(c\) contributes
\([T,\Delta_c|_{S=0}]=0\).  Thus no outside-letter term remains,
and \eqref{eq:special-fb} gives
\begin{equation}\label{eq:special-transverse}
 \begin{split}
 -[T,\psi(c,a)]={}&k(\ad_a^n-\ad_c^n)T\\
 &+[a,\partial_2\psi(b,a)T]
       -[c,\partial_2\psi(b,c)T].
 \end{split}
\end{equation}

Now define
\(G(a,b)=[a,\psi(b,a)]+[c,\psi(b,c)]\), where \(c=-a-b\).
Hold \(b\) fixed and vary \(a\) in direction \(T\), so that \(c\)
varies in direction \(-T\).  The full product rule is
\begin{align*}
 d_aG(T)={}&[T,\psi(b,a)-\psi(b,c)]\\
           &+[a,\partial_2\psi(b,a)T]
              -[c,\partial_2\psi(b,c)T].
\end{align*}
The three-term identity and antisymmetry give
\(\psi(b,a)-\psi(b,c)=\psi(c,a)\).  Comparing with
\eqref{eq:special-transverse}, we conclude that
\begin{equation}\label{eq:special-dg}
 d_aG(T)=-k(\ad_a^n-\ad_c^n)T.
\end{equation}

\emph{Vanishing of the scalar obstruction.}
Equation~\eqref{eq:special-dg} expresses the derivative of \(G\) in
terms of the single scalar \(k\).  To show that this scalar is zero,
introduce two independent constant directions \(U,V\).  The
coefficients of their two commuting scalar parameters in \(G\)
are unchanged when the differentiations are interchanged.  After
forming these mixed derivatives, specialize
\[
 a=a_0,\qquad b=-a_0,\qquad c=0,\qquad U=a_0,\qquad V=T,
\]
where \(a_0,T\) are free letters.  The derivative of \(\ad_c^n\)
vanishes there, since every summand contains \(n-1>0\) copies of
\(\ad_c\).  The two derivatives of the other term are
\begin{align*}
 d_U(\ad_a^nV)\big|&=n\ad_{a_0}^nT,\\
 d_V(\ad_a^nU)\big|
   &=\sum_{j=0}^{n-1}
       \ad_{a_0}^{j}\ad_T\ad_{a_0}^{n-1-j}a_0
     =-\ad_{a_0}^{n}T.
\end{align*}
Only \(j=n-1\) survives in the second sum.  Equality of the mixed
derivatives in \eqref{eq:special-dg} implies
\[
 (n+1)k\ad_{a_0}^{n}T=0.
\]
The word \(a_0^nT\) has coefficient one in
\(\ad_{a_0}^{n}T\).  Since the calculation is over \(\Q\), it
follows that \(k=0\).  This proves
\eqref{eq:special-linearization}.

Equation~\eqref{eq:special-dg} now says that every directional
derivative of \(G\) in \(a\) vanishes.  Equivalently, setting the
direction equal to \(a\) multiplies its component with \(j\)
occurrences of \(a\) by \(j\).  In characteristic zero every such
component with \(j>0\) is zero.  Thus \(G\) is independent of
\(a\), and evaluation at \(a=0,c=-b\) gives \(G=0\).  This is
\eqref{eq:special}.
\end{proof}

\begin{proposition}[Saturated reduction]\label{prop:saturated-reduction}
Let
\[
 K_n=\Lie_R\langle x,y\rangle_n,
 \qquad L_n=W_n\cap K_n.
\]
Then \(L_n\) is a saturated finite free \(R\)-submodule of \(K_n\),
and its reduction embeds in \(K_n/2K_n\).  Writing its image as
\(\overline L_n\), one has
\begin{equation}\label{eq:saturated-dimension}
 \dim_{\F}\overline L_n=\dim_{\Q}W_n.
\end{equation}
More generally, for any rational subspace \(V_n\subset W_n\), the
lattice \(M_n=V_n\cap K_n\) has an embedded reduction
\(\overline M_n\subset\overline L_n\) of dimension
\(\dim_{\Q}V_n\).

Every element of \(\overline L_n\) satisfies the reductions of the
three defining identities.  If \(n>2\), it also satisfies
\eqref{eq:special} and \eqref{eq:special-linearization} after
reduction, together with all formal coefficient consequences of
these identities.  In particular, putting \(z=x+y\),
\begin{equation}\label{eq:special-mod-two}
 [x,\overline\psi(x,y)+\overline\psi(z,y)]
       +[y,\overline\psi(z,y)]=0
 \qquad(\overline\psi\in\overline L_n).
\end{equation}
\end{proposition}
\begin{proof}
An intersection of a rational subspace with a finite free
\(R\)-module is saturated: if \(v\in K_n\) and \(2v\in L_n\),
then \(v\in W_n\), hence \(v\in L_n\).  Thus
\(L_n\cap2K_n=2L_n\), which gives the asserted injection after
reduction.  It is a finite free module because \(R\) is a principal
ideal domain.  Clearing denominators in a rational basis of
\(W_n\) shows that \(\Q L_n=W_n\), so its rank and its reduction
dimension are \(\dim_{\Q}W_n\).  Exactly the same argument applies
to \(M_n\); embedding both reductions in \(K_n/2K_n\) proves their
stated inclusion.

For a representative \(\psi\in L_n\), each defining identity is
an expression with coefficients in \(R\).  Its target is either a
free Lie lattice or the braid Lie algebra over \(R\).  These targets
embed in their rationalizations by Lemmas~\ref{lem:lie-lattice} and
\ref{lem:free-fiber}.  An expression which vanishes over \(\Q\)
therefore already vanishes over \(R\), and can be reduced modulo
two.  The identical argument applies to the special identity and
its linearization, which were established over \(\Q\) in
Lemma~\ref{lem:special}.  Adjoining free letters and commuting
parameters preserves the word bases; extracting a parameter
coefficient involves no division and commutes with reduction.
This justifies the assertion about formal coefficient consequences.

Finally, reduce \eqref{eq:special}, take \(a=x,b=y,c=x+y\), and
use the reduced antisymmetry to replace \(\overline\psi(y,x)\)
and \(\overline\psi(y,x+y)\) by
\(\overline\psi(x,y)\) and \(\overline\psi(x+y,y)\).
Expanding the outside bracket with \(x+y\) gives
\eqref{eq:special-mod-two}.
\end{proof}

In the characteristic-two arguments that follow, the permitted
polynomials are these reductions of rational solutions.  No
identification of \(\overline L_n\) with the full solution space of
equations written over \(\F\) is asserted or needed.  In particular,
the characteristic-zero mixed-derivative argument proving
\((n+1)k=0\) has already been completed before reduction.
